% © 2025 Ing. David Jaroš — CC BY-NC-ND 4.0
%
% This work is licensed under a Creative Commons Attribution-NonCommercial-NoDerivatives
% 4.0 International License (CC BY-NC-ND 4.0).
%
% License History: Earlier drafts (up to v0.3) were released under CC BY 4.0.
% From v0.4 onward, all material is released under CC BY-NC-ND 4.0 to protect
% the integrity of the theoretical work during ongoing academic development.

% =====================================================================
% TASK 1 — Hermitian winding-vacuum audit and Schwarzschild solution
% Target: canonical/geometry/biquaternionic_vacuum_solutions.tex
% Status: §1 [DERIVED — CORRECTED], §2 [SKETCH], §3 [DERIVED]
% Layer: [L1] — biquaternionic geometry only (no SM input)
% =====================================================================

\section{Hermitian Winding-Vacuum Audit and Schwarzschild Solution}
\label{sec:biq_vacuum_solutions}

\noindent\textbf{Notation:} Throughout, $\varphi$ denotes the phase-frame parameter in
$P_\varphi[\mathcal{G}_{\mu\nu}]$ (observer choice), and $\psi$ is the imaginary part of complex
time $\tau = t + i\psi$ (dynamical field).  These are distinct objects and must not be conflated.

\noindent\textbf{Layer:} [L1] — geometry only; no Standard Model input required.

\noindent\textbf{Key references:}
\begin{itemize}
  \item Biquaternionic metric postulate: \texttt{canonical/geometry/biquaternion\_metric.tex}
  \item Tetrad definition: \texttt{canonical/geometry/biquaternion\_tetrad.tex}
  \item Phase-frame projection: \texttt{canonical/geometry/phase\_projection.tex}
  \item $\varphi$ gauge vs.\ physical: \texttt{canonical/geometry/phi\_gauge\_vs\_physical.tex}
  \item Supporting script: \texttt{tools/compute\_h\_munu\_vacuum.py}
\end{itemize}

---

\subsection{Motivation}
\label{sec:biq_vacuum_motivation}

From \texttt{canonical/geometry/phi\_gauge\_vs\_physical.tex} (Gap 3, [DERIVED]):
\begin{equation}
  \frac{\partial\alpha}{\partial\varphi}\bigg|_{\varphi=0} = 2\rho\, r\,\alpha(0),
  \qquad r := \frac{|\mathcal{A}^I_\mu|}{|\mathcal{A}^R_\mu|}.
  \label{eq:dalpha_dphi_recap}
\end{equation}
The diagnostic can establish a physical phase modulus only after a nonzero
imaginary sector has been derived from a well-defined UBT observable.  The
Hermitian metric channel used below does not provide such a sector: the
quantity $\mathrm{Sc}(E_\psi E_\psi^\dagger)$ is real for every field, not
only for a single winding mode.  This section records the corrected no-go and
keeps the separate Schwarzschild construction.

---

\subsection{Single-Mode Winding: Dead End for $h_{\mu\nu}$}
\label{sec:single_mode_dead_end}

\subsubsection{Ansatz}

Consider the simplest non-trivial winding background:
\begin{equation}
  \Theta(x,\psi) = \Theta_0 \cdot e^{i\psi/R_\psi},
  \qquad \Theta_0 \in \mathbb{B},\; R_\psi > 0.
  \label{eq:single_mode_ansatz}
  \quad \textbf{[POSTULATE — single-mode ansatz, winding $n=1$]}
\end{equation}

\subsubsection{Computation}

The tetrad component [DERIVED from Eq.\,\eqref{eq:single_mode_ansatz}]:
\begin{equation}
  E_\psi = \partial_\psi\Theta = \frac{i}{R_\psi}\,\Theta_0\,e^{i\psi/R_\psi}.
\end{equation}

The $\psi\psi$-component of the biquaternionic metric
$\mathcal{G}_{\mu\nu} = \mathrm{Sc}(E_\mu E_\nu^\dagger)$ [POSTULATE — tetrad formula]:
\begin{align}
  \mathcal{G}_{\psi\psi}
  &= \mathrm{Sc}(E_\psi\,E_\psi^\dagger) \notag\\
  &= \mathrm{Sc}\!\left(\frac{i}{R_\psi}\Theta_0\,e^{i\psi/R_\psi}
       \cdot \frac{-i}{R_\psi}\Theta_0^\dagger\,e^{-i\psi/R_\psi}\right) \notag\\
  &= \frac{1}{R_\psi^2}\,\mathrm{Sc}(\Theta_0\,\Theta_0^\dagger)
   = \frac{|\Theta_0|^2}{R_\psi^2} \in \mathbb{R}.
  \label{eq:single_mode_G_psipsi}
  \quad \textbf{[DERIVED]}
\end{align}

Here $|\Theta_0|^2 = \mathrm{Sc}(\Theta_0\,\Theta_0^\dagger)$ is the biquaternionic
norm-squared, which is real and non-negative (standard result).

\subsubsection{Conclusion}

\begin{equation}
  h_{\psi\psi} = \mathrm{Im}(\mathcal{G}_{\psi\psi}) = 0
  \quad \text{for the single-mode ansatz.}
  \label{eq:single_mode_dead_end}
  \quad \textbf{[DERIVED — DEAD END for $h_{\mu\nu}$]}
\end{equation}

A single winding mode cannot generate an imaginary metric component.
This is recorded explicitly to guide future attempts.

---

\subsection{Two-Mode Hermitian Channel: Correction}
\label{sec:two_mode_vacuum}

\noindent\textbf{Correction (31 July 2026).}
The previous version claimed that two winding modes generate an imaginary
component of $\mathcal G_{\psi\psi}=\mathrm{Sc}(E_\psi E_\psi^\dagger)$.
That claim was algebraically false.  A Hermitian norm is real, and the two
cross terms are complex conjugates.  Hence $h_{\psi\psi}=0$ identically in
this channel.  The earlier conclusion that this example proves a physical
$\varphi$ parameter is withdrawn.

\subsubsection{Ansatz}

Keep the two-mode profile
\begin{equation}
  \Theta(x,\psi)=\Theta_0e^{i\psi/R_\psi}
                 +\Theta_1e^{2i\psi/R_\psi},
  \label{eq:two_mode_ansatz}
  \qquad \Theta_0,\Theta_1\in\mathbb B.
\end{equation}
Then
\begin{equation}
  E_\psi=\frac{i}{R_\psi}\Theta_0e^{i\psi/R_\psi}
        +\frac{2i}{R_\psi}\Theta_1e^{2i\psi/R_\psi}.
  \label{eq:two_mode_E_psi}
\end{equation}
Let
\[
 N_0:=\mathrm{Sc}(\Theta_0\Theta_0^\dagger),\qquad
 N_1:=\mathrm{Sc}(\Theta_1\Theta_1^\dagger),\qquad
 s_{01}:=\mathrm{Sc}(\Theta_0\Theta_1^\dagger).
\]
The two cross terms are
\begin{equation}
 \frac{2}{R_\psi^2}\left(
 e^{-i\psi/R_\psi}s_{01}
 +e^{i\psi/R_\psi}\overline{s_{01}}
 \right)
 =\frac{4}{R_\psi^2}\left[
 \cos\!\left(\frac{\psi}{R_\psi}\right)\mathrm{Re}(s_{01})
 +\sin\!\left(\frac{\psi}{R_\psi}\right)\mathrm{Im}(s_{01})
 \right]\in\mathbb R.
 \label{eq:cross_sum}
\end{equation}
Therefore the complete Hermitian metric component is
\begin{equation}
 \boxed{
 \mathcal G_{\psi\psi}
 =\frac1{R_\psi^2}\left[
 N_0+4N_1
 +4\cos\!\left(\frac{\psi}{R_\psi}\right)\mathrm{Re}(s_{01})
 +4\sin\!\left(\frac{\psi}{R_\psi}\right)\mathrm{Im}(s_{01})
 \right]\in\mathbb R,}
 \label{eq:cross_sum_expanded}
\end{equation}
and hence
\begin{equation}
 \boxed{h_{\psi\psi}=\mathrm{Im}(\mathcal G_{\psi\psi})=0.}
 \label{eq:h_psipsi_result}
 \qquad\textbf{[DERIVED --- CORRECTED]}
\end{equation}
The statement is stronger than a two-mode cancellation.  For any
biquaternionic profile $E_\psi$, the scalar Hermitian norm
$\mathrm{Sc}(E_\psi E_\psi^\dagger)$ is real.  Adding further winding modes
cannot generate an imaginary component through this same metric definition.

\subsubsection{Canonical numerical audit}

For the former example
\begin{equation}
 \Theta_0=(1+i_c)\mathbf 1,\qquad
 \Theta_1=\mathbf 1+i_c\mathbf j,
 \label{eq:canonical_theta_choice}
\end{equation}
we have $N_0=N_1=2$ and $s_{01}=1+i_c$, so
\begin{equation}
 \mathcal G_{\psi\psi}
 =\frac{10+4\cos(\psi/R_\psi)+4\sin(\psi/R_\psi)}{R_\psi^2},
 \qquad h_{\psi\psi}=0.
 \label{eq:h_psipsi_canonical}
\end{equation}
The supporting script
\texttt{tools/compute\_h\_munu\_vacuum.py} now verifies this identity directly.

---

\subsection{Gauge-Potential Sketch and $\varphi$ Status}
\label{sec:r_and_dalpha_dphi}

The exploratory expression
\begin{equation}
 \mathcal A_\psi\approx
 \frac{\mathrm{Sc}(\Theta^\dagger\partial_\psi\Theta)}{|\Theta|^2}
 \label{eq:gauge_potential_sketch}
 \qquad\textbf{[SKETCH]}
\end{equation}
can still return a complex diagnostic.  The current script reports, for the
canonical example and the normalization used there,
\begin{align}
 r&=\frac{|\mathcal A^I_\psi|}{|\mathcal A^R_\psi|}\approx4.66,
 \label{eq:r_numerical}\\
 \rho&\approx0,\qquad
 \frac{\partial\alpha}{\partial\varphi}\bigg|_{\varphi=0}
 =2\rho r\alpha(0)\approx0.
 \label{eq:dalpha_dphi_numerical}
\end{align}
This number is not a derivation of a physical imaginary metric sector.  The
potential formula is only a sketch, its relation to the canonical QED sector
has not been proved, and $r\neq0$ by itself does not imply
$\partial\alpha/\partial\varphi\neq0$.  Consequently the example does not
establish that $\varphi$ is physical.

\subsection{Corrected Status}
\label{sec:biq_vacuum_summary}

\begin{center}
\begin{tabular}{p{0.30\linewidth}p{0.18\linewidth}p{0.42\linewidth}}
\hline
Item & Status & Result \\
\hline
Single-mode Hermitian channel & Derived & $h_{\psi\psi}=0$. \\
Two-mode Hermitian channel & Derived, corrected & $h_{\psi\psi}=0$; conjugate cross terms are real. \\
Gauge diagnostic $r\approx4.66$ & Sketch only & Does not prove $h_{\mu\nu}\neq0$ or physical $\varphi$. \\
Physical $\varphi$ vacuum & Open & Requires a separately derived non-Hermitian observable or another canonical imaginary sector with $\rho r\neq0$. \\
\hline
\end{tabular}
\end{center}

The repository status must therefore be downgraded from ``proved physical'' to
``open''.  No claim about the later Schwarzschild solution depends on the
withdrawn two-mode argument.

% =====================================================================
% §3 — Full Schwarzschild Vacuum Solution [PROMOTED FROM RESEARCH_TRACKS v57]
% Source: research_tracks/research/schwarzschild_from_theta.tex §5
% Status: Proved [L1]
% =====================================================================

\section{Full Schwarzschild Vacuum Solution}
\label{sec:schwarzschild_vacuum}

\noindent\textbf{Status:} \textbf{Proved [L1]} — algebraic computation plus complex-time
structure; verified numerically.

\noindent\textbf{Source:}
\texttt{research\_tracks/research/schwarzschild\_from\_theta.tex} (full derivation).
Numerical verification: \texttt{tools/verify\_schwarzschild\_theta.py} (relative error
$< 10^{-8}$ at $r/M = 2,5,10,50,100$).

---

\subsection{Ansatz}

\begin{equation}
  \boxed{\Theta_0(r,\tau) = e^{i\alpha(r)}\cdot\bigl[f(r)\cdot\mathbf{1}
    + g(r)\cdot\boldsymbol{e}_r\bigr],
  \quad \tau = t + i\psi,\quad |\Theta_0| = 1.}
  \label{eq:schwarzschild_ansatz}
  \quad\textbf{[POSTULATE — spherically symmetric ansatz]}
\end{equation}
Here $\boldsymbol{e}_r = (x_i/r)\boldsymbol{e}_i$ is the radial unit biquaternion,
$f(r),g(r) \in \mathbb{R}$, and $\alpha(r) = \Phi(r)$ is fixed by the matching
condition (see below).

---

\subsection{Spatial Metric: Phase Cancellation Identity}

\begin{theorem}[Phase cancellation, \textbf{[L1]}]
\label{thm:phase_cancel}
For any radial phase $e^{i\alpha(r)}$ and spatial indices $i,j$:
\begin{equation}
  g_{ij} = \mathrm{Sc}\bigl((\partial_i\Theta_0)^\dagger\,\partial_j\Theta_0\bigr)
          = \mathrm{Sc}\bigl(\widetilde{E}_i^\dagger\,\widetilde{E}_j\bigr),
  \label{eq:phase_cancel}
\end{equation}
where $\widetilde{E}_i := \partial_i[f\mathbf{1}+g\boldsymbol{e}_r]$ is the spatial
derivative without the phase factor.  The phase cancels because
$\mathrm{Sc}(e^{-i\alpha}Ae^{i\alpha}) = \mathrm{Sc}(A)$ for all
$A \in \mathcal{B}$ (Sc is $\mathbb{C}$-linear and commutes with conjugation).
\end{theorem}

\noindent\textbf{Spatial metric result [DERIVED]:}
The ODE system
\begin{equation}
  f'(r)^2 + g'(r)^2 = \Psi(r)^4,\qquad
  g(r) = r\,\Psi(r)^2,\qquad
  \Psi(r) = 1 + \frac{M}{2r},
  \label{eq:schwarzschild_ode}
\end{equation}
gives $g_{ij}[\Theta_0] = \Psi(r)^4\,\delta_{ij}$ — the isotropic Schwarzschild
spatial metric.

---

\subsection{Temporal Metric: Complex Time Structure}

From the complex-time signature theorem
(\texttt{canonical/gr\_closure/step3\_signature\_theorem.tex}), the
$\psi$-derivative gives the temporal metric component:
\begin{equation}
  g_{tt} = -|\partial_\psi\Theta_0|^2.
  \label{eq:g_tt_psi}
\end{equation}

With the ansatz $\Theta_0 = e^{i\alpha(r)}\cdot X_0(r)$, $|X_0|=1$, the
$\psi$-derivative (static ansatz: $X_0$ independent of $\psi$) gives:
\begin{equation}
  \partial_\psi\Theta_0 = i\,\alpha(r)\,\Theta_0,
  \qquad\Longrightarrow\qquad
  g_{tt} = -\alpha(r)^2.
  \label{eq:g_tt_alpha}
\end{equation}

\noindent\textbf{Matching condition:} Setting $g_{tt} = -\Phi(r)^2$ requires:
\begin{equation}
  \alpha(r) = \Phi(r) = \frac{1 - M/2r}{1 + M/2r}.
  \label{eq:alpha_matching}
  \quad\textbf{[DERIVED]}
\end{equation}

---

\subsection{Main Theorem}

\begin{theorem}[Full Schwarzschild metric from $\Theta_0$, \textbf{Proved [L1]}]
\label{thm:schwarzschild_full}
The ansatz \eqref{eq:schwarzschild_ansatz} with $\alpha(r) = \Phi(r)$,
$g(r) = r\Psi(r)^2$, $f'(r) = \Psi(r)\sqrt{2M/r}$, and complex time
$\tau = t + i\psi$ gives the full Schwarzschild metric in isotropic coordinates:
\begin{equation}
  \boxed{g_{tt} = -\Phi(r)^2,\qquad
  g_{ij} = \Psi(r)^4\,\delta_{ij},}
  \quad \Phi = \frac{1-M/2r}{1+M/2r},\quad \Psi = 1+\frac{M}{2r}.
  \label{eq:schwarzschild_full}
\end{equation}
\end{theorem}

\begin{proof}[Proof sketch]
Spatial sector: Theorem~\ref{thm:phase_cancel} reduces to the phase-free
computation; the ODE \eqref{eq:schwarzschild_ode} gives $g_{ij}=\Psi^4\delta_{ij}$
(proved in \texttt{schwarzschild\_from\_theta.tex §3–4}).
Temporal sector: Eq.~\eqref{eq:g_tt_alpha} with matching condition
\eqref{eq:alpha_matching} gives $g_{tt}=-\Phi^2$.
The two sectors are independent because phase cancellation
(Theorem~\ref{thm:phase_cancel}) decouples them.
\end{proof}

\begin{tcolorbox}[colback=green!5!white,colframe=green!75!black,
  title={Self-Containedness Check}]
\begin{itemize}
  \item Uses only: signature theorem (\texttt{step3\_signature\_theorem.tex}),
    scalar inner product algebra, and the spherically symmetric ansatz.
  \item No free parameters: $\alpha(r) = \Phi(r)$ is fixed by matching
    $g_{tt} = -\Phi^2$.
  \item Numerically verified: \texttt{tools/verify\_schwarzschild\_theta.py}
    gives relative error $< 10^{-8}$ at $r/M \in \{2,5,10,50,100\}$.
  \item Layer: \textbf{[L1]} — biquaternionic geometry + complex time; no SM input.
\end{itemize}
\end{tcolorbox}

---

% =====================================================================
% TASK 3 — Independent Fixation of R_ψ
% Status: §2 [KALIBROVÁNO — no independent topological derivation found]
% Layer: [L1] — biquaternionic geometry only
% =====================================================================

\section{Independent Fixation of $R_\psi$: Topological Analysis}
\label{sec:Rpsi_fixation}

\noindent\textbf{Status:} [KALIBROVÁNO — nutná pro konzistenci; nederivovatele z čisté
geometrie v současné teorii]

\noindent\textbf{Problem.}
The radius $R_\psi$ of the imaginary-time circle $S^1_\psi$ is currently identified as
\begin{equation}
  R_\psi = \frac{\hbar}{m_e c}
  \label{eq:Rpsi-current}
\end{equation}
[CALIBRATED — chosen so that $\Lambda_\psi = \hbar c/R_\psi = m_e c^2$].
This identification collapses two independent scales ($R_\psi$ and $m_e$) into one.
This section tests whether a purely topological condition can fix $R_\psi$
independently of $m_e$.

---

\subsection{Candidate 1: Self-Duality of the Torus}

The torus $S^1_t\times S^1_\psi$ is self-dual (invariant under modular transformation
$\tau\to-1/\tau$) if and only if $R_t = R_\psi$.  From $\alpha = R_t/R_\psi$, this
requires $\alpha=1$.  \textbf{[DEAD END — $\alpha=1$ is unphysical.]}

---

\subsection{Candidate 2: Winding Consistency on $S^1_\psi$}
\label{sec:winding-consistency}

Single-valuedness of the charged $\Theta$-field under $n^\star$ complete traversals
of the $\psi$-circle requires the accumulated phase to be a multiple of $2\pi$:
\begin{equation}
  n^\star \cdot \frac{2\pi R_\psi \, m_e c}{\hbar} = 2\pi k\,,\quad k\in\mathbb{Z}\,.
  \label{eq:winding-consistency}
\end{equation}
This gives
\begin{equation}
  R_\psi = \frac{k\,\hbar}{n^\star\, m_e c}\,.
  \label{eq:Rpsi-winding}
\end{equation}
The two simplest cases are:

\begin{center}
\begin{tabular}{|c|l|l|l|}
\hline
$k$ & $R_\psi$ & $\Lambda_\psi = \hbar c/R_\psi$ & Consequence for $\alpha$ \\
\hline
$1$ & $\hbar/(n^\star m_e c)=\hbar/(137\,m_e c)$ & $n^\star m_e c^2 = 137\,m_e c^2$ &
  $\alpha = m_0/(n^{\star\,2} m_e) \approx 1/137^2$ \ [too small] \\
$n^\star$ & $n^\star\hbar/(n^\star m_e c)=\hbar/(m_e c)$ & $m_e c^2$ &
  recovers Eq.~\eqref{eq:Rpsi-current} [original identification] \\
\hline
\end{tabular}
\end{center}

Neither $k=1$ nor $k=n^\star$ provides a genuinely new, independent fixing of $R_\psi$:
the former gives the wrong order of magnitude for $\alpha$; the latter reproduces
the calibrated identification.

\textbf{[DEAD END for $k=1$; REPRODUCES CALIBRATION for $k=n^\star$.]}

---

\subsection{Candidate 3: Modularity / Theta-Function Consistency}

The UBT partition function involves theta functions on $S^1_\psi$.  For the partition
function to transform covariantly under $S\in SL(2,\mathbb{Z})$ (modular group), one
needs $R_\psi^2 \in \mathbb{Q}$ in natural units.  This constrains $R_\psi$ to a
dense set, not a unique value — it does not pin down $R_\psi$ without additional
input.  \textbf{[DEAD END — modularity alone does not fix $R_\psi$]}

---

\subsection{Honest Summary: Status of $R_\psi$}

\begin{center}
\fbox{\parbox{0.88\linewidth}{%
\textbf{Conclusion [KALIBROVÁNO]:}\\
The identification $R_\psi = \hbar/(m_e c)$ is required for internal consistency
(it ensures $\Lambda_\psi = m_e$, which closes the fermion mass spectrum), but
no purely topological or geometric condition has been found that derives this value
\emph{independently} of $m_e$.\\[4pt]
\textbf{What would change this:}\\
A derivation of $R_\psi$ from the UBT action alone (e.g., from the spectrum of
$\mathcal{L}_\text{geom}$ on $S^1_\psi$) without referring to any lepton mass.
}}
\end{center}

\noindent See \texttt{docs/ALPHA\_FROM\_ME\_ANALYSIS.md} §6.2 for numerical details
and the impact on $\alpha$ predictions.

\noindent See \texttt{docs/FITTED\_PARAMETERS.md} for the updated parameter status.
